\section*{Chapter \ref{ch:b1:simple-regression} --- Simple Regression}
% Numbers recomputed in code/bachelor-1/32-simple-regression/exercises/ (solutions.py,
% exo_07.py; tested in tests/test_ch32.py) and in the chapter's captured outputs.

\begin{solution}{exo:b1:simple-regression:five}
$\bar x = 12$, $\bar y = 9.84$; $\sum (x_i - \bar x)(y_i - \bar y) = 4.4$ and $\sum (x_i - \bar x)^2 = 40$, so $\hat\beta_1 = 0.11$ and
$\hat\beta_0 = 9.84 - 0.11 \times 12 = 8.52$.
\end{solution}

\begin{solution}{exo:b1:simple-regression:read}
Slope $0.0885$, \omterm{def:b1:sampling-and-inference:standard-error}{standard error} $0.005$ ($0.00487$ unrounded), $R^2 = 0.1304$.
\end{solution}

\begin{solution}{exo:b1:simple-regression:interpret}
One more year of experience goes with wages about $6\%$ higher ($100(e^{0.06} - 1) = 6.18\%$
exactly), on average across workers; it is a description, not yet an effect.
\end{solution}

\begin{solution}{exo:b1:simple-regression:predict}
$8.8172 + 16 \times 0.0885 = 10.233$, which corresponds to $e^{10.233} \approx 27{,}800$ dollars --- a
typical (geometric-mean) value, below the mean earnings of such men.
\end{solution}

\begin{solution}{exo:b1:simple-regression:residuals}
$\sum_i \hat y_i\hat u_i = \hat\beta_0\sum_i \hat u_i + \hat\beta_1\sum_i x_i\hat u_i = 0$ by \cref{prop:b1:simple-regression:residuals}, and the \omterm{def:b1:simple-regression:residual}{residuals}
have mean zero, so their sample covariance with $\hat y$ is zero.
\end{solution}

\begin{solution}{exo:b1:simple-regression:ages}
Under 30: 953 men, slope $0.117$ (\omterm{def:b1:sampling-and-inference:standard-error}{standard error} $0.0097$). Over 40: 678 men, slope
$0.0939$ ($0.0074$). The difference, $0.023$, has a \omterm{def:b1:sampling-and-inference:standard-error}{standard error} of about
$\sqrt{0.0097^2 + 0.0074^2} = 0.012$: about two \omterm{def:b1:sampling-and-inference:standard-error}{standard errors}, suggestive but not
conclusive.
\end{solution}

\begin{solution}{exo:b1:simple-regression:energy}
183 countries; slope $1.247$, \omterm{def:b1:sampling-and-inference:standard-error}{standard error} $0.033$: $1\%$ more income per person goes
with $1.25\%$ more energy use per person, as \cref{ch:b1:comparative-statics} found.
\end{solution}

\begin{solution}{exo:b1:simple-regression:direction}
Upward: wealth raises earnings and is positively correlated with
schooling, so the slope of earnings on schooling picks up part of
wealth's effect.
\end{solution}

\begin{solution}{exo:b1:simple-regression:test}
$t = (0.0885 - 0.10)/0.00487 = -2.36$, beyond $-1.96$: the slope is significantly below
$0.10$ at $5\%$.
\end{solution}

\begin{solution}{exo:b1:simple-regression:derive}
Setting the derivatives of $\sum_i (y_i - b_0 - b_1x_i)^2$ to zero gives $\sum_i \hat u_i = 0$ and $\sum_i x_i\hat u_i = 0$.
The first gives $\hat\beta_0 = \bar y - \hat\beta_1\bar x$; substituting into the second and subtracting
$\bar x\sum_i \hat u_i = 0$ gives $\sum_i (x_i - \bar x)(y_i - \bar y - \hat\beta_1(x_i - \bar x)) = 0$, hence $\hat\beta_1$. Unbiasedness:
write $y_i = \beta_0 + \beta_1x_i + u_i$, so $\hat\beta_1 = \beta_1 + \sum_i (x_i - \bar x)u_i/\sum_i (x_i - \bar x)^2$; given the $x$'s, each
$\E[u_i \mid x] = 0$, so the expectation is $\beta_1$.
\end{solution}

\begin{solution}{exo:b1:simple-regression:product}
The two slopes are $s_{xy}/s_x^2$ and $s_{xy}/s_y^2$; their product is $s_{xy}^2/(s_x^2s_y^2) = r^2$. With $r^2 < 1$ the
two cannot both be at least one when the variances are similar; here
$0.8368 \times 0.6394 = 0.5350 = 0.7315^2$. Earnings in each year are partly transitory, so
each year's earnings predict the other's with a slope below one.
\end{solution}

\begin{solution}{exo:b1:simple-regression:refute}
Twins or siblings with different schooling but the same earnings; or a
natural experiment --- a change in school-leaving rules that raised schooling
for some cohorts --- with no matching rise in their earnings; or a slope that
falls to near zero once measures of ability and family background enter.
\end{solution}

\begin{solution}{pb:b1:simple-regression:brenland}
\textbf{1.} 1,922.

\textbf{2.} Mean $12.63$ years, standard deviation $2.59$.

\textbf{3.} Slope $0.0580$, intercept $9.3956$.

\textbf{4.} The sum is $1.3 \times 10^{-11}$: zero up to rounding error.

\textbf{5.} $10.091$ and $10.323$.

\textbf{6.} $e^{4 \times 0.058} = 1.261$: $26\%$ higher.

\textbf{7.} 1,841 dalens per year of schooling.

\textbf{8.} A year more of the head's schooling goes with about $5.8\%$ more
household income.

\textbf{9.} $100(e^{0.058} - 1) = 5.97\%$.

\textbf{10.} $0.108$.

\textbf{11.} No: it means schooling leaves most of the variation unexplained;
the slope is precisely estimated (item 12).

\textbf{12.} $0.0038$; $t = 15.3$.

\textbf{13.} $0.0505$ to $0.0654$.

\textbf{14.} $t = (0.058 - 0.10)/0.0038 = -11.1$: rejected.

\textbf{15.} $0.1888$.

\textbf{16.} $0.0628$.

\textbf{17.} Schooling $0.0453$, ability $0.0927$.

\textbf{18.} $0.0628 - 0.0453 = 0.0175 = 0.0927 \times 0.1888$.

\textbf{19.} By $0.0128$ (from $0.05$ to $0.0628$).

\textbf{20.} Ability raises income and goes with more schooling, so the short
regression credits schooling with part of ability's effect.

\textbf{21.} Yes, if anything that raises income is also correlated with
schooling; the survey has no ability measure to check.

\textbf{22.} $0.0591$.

\textbf{23.} $-0.0034$: $0.34\%$ less income per year of age, schooling held fixed.

\textbf{24.} $0.049$: age is almost uncorrelated with schooling, so leaving it out
biases the slope very little.

\textbf{25.} When the variable both affects income and is correlated with
schooling, like ability in Part III; this is what multiple regression in
the second university year handles.

\textbf{26.} The Brenland return to a year of the head's schooling is $5.8\%$
(\omterm{def:b1:sampling-and-inference:standard-error}{standard error} $0.4$ point) in the survey; if it carried an ability bias of
the size simulated in Part III (1.3 points), the causal return would be
about $4.5\%$.
\end{solution}
